\section{Kiểm thử và phạm vi đóng góp}

\subsection{Kiểm thử tích hợp}

Các nhóm kiểm thử tập trung vào hành vi có thể quan sát:
\begin{itemize}
    \item UCS: đối chiếu chi phí với BFS, replay hợp lệ, không làm đổi
    trạng thái đầu vào, xử lý chu trình và đường đi có chi phí cải thiện.
    \item Benchmark: lời giải của cả bốn thuật toán, số node ở các
    trường hợp nhỏ, trạng thái đã ở đích/vô nghiệm, giới hạn và dọn bộ đo
    khi lỗi xảy ra.
    \item UI: worker thật dùng \texttt{spawn}, chạy và hủy tìm kiếm,
    replay, đổi khối con, xuất CSV và không làm thay đổi tiến độ chiến
    dịch khi dùng solver.
    \item Visualization: phép tổng hợp và tỷ lệ theo cùng map; phát
    hiện dữ liệu trùng, thiếu, trộn run, sai chi phí hoặc replay không
    hợp lệ; không thay số liệu thiếu bằng 0.
\end{itemize}

Tại merge commit \texttt{8f1f183}, bộ kiểm thử ghi nhận 237 tests đạt.
Sau khi cập nhật \texttt{main} tới \texttt{a89aad2}, bộ kiểm thử hiện có
ghi nhận 228 tests đạt và 8 kiểm tra visualization đạt. Hai con số
thuộc hai phiên bản có tập test khác nhau; bản cập nhật đã bỏ file
test DFS riêng và bổ sung kiểm tra maps mới, nên không diễn giải sự
thay đổi số lượng thành mức tăng hoặc giảm chất lượng phần mềm.
Các lệnh kiểm tra tương ứng từ thư mục gốc:

\begin{lstlisting}[language=bash]
python -m pytest -q
python -m unittest visualization.checks -v
\end{lstlisting}

Các kết quả trên xác nhận những trường hợp được kiểm thử, không thay
thế chứng minh tính đúng của thuật toán hoặc bảo đảm không còn lỗi
trên mọi đầu vào. Phép đo hiệu năng trong Mục~\ref{sec:benchmark}
vẫn dùng CSV lịch sử, không lấy thời gian chạy test làm benchmark.

\subsection{Phạm vi phần việc được trình bày}

Phần đóng góp của Nguyễn Thành Bảo được mô tả trong các mục trên gồm:
cài đặt UCS; tích hợp chế độ chạy thuật toán và replay vào UI; chuẩn
hóa phép đo benchmark; xây dựng CSV 20 maps và schema; tạo notebook
với bảy biểu đồ; xử lý tương thích và kiểm thử khi merge với code nhóm.

Bộ luật chơi, giao diện nền, BFS, DFS và heuristic A* là các thành
phần chung của nhóm được sử dụng hoặc điều chỉnh để tích hợp. Bộ 15
maps mới trên \texttt{main} là phần bổ sung của nhóm, không được nhận
là phần tự xây dựng trong phạm vi công việc này. Mức độ hoàn thành
và phân công cuối cùng cần được các thành viên thống nhất.

Trong quá trình thực hiện có sử dụng trợ lý AI Codex để hỗ trợ soạn
và chỉnh sửa code, tạo kiểm thử, xử lý tích hợp Git và tổng hợp nội
dung báo cáo. Bản thảo này được đối chiếu với mã nguồn, CSV và kết
quả kiểm thử trong repository; người thực hiện cần rà soát nội dung
và tuân theo quy định khai báo hỗ trợ AI của môn học trước khi nộp.
